\chapter{การทบทวนวรรณกรรมเชิงธีมและการวิเคราะห์ช่องว่างงานวิจัย}

การทบทวนวรรณกรรมในงานวิจัยระดับสูง ไม่ใช่เพียงการนำบทคัดย่อของบทความหลายเรื่องมาเรียงต่อกัน แต่คือกระบวนการวิเคราะห์เชิงสังเคราะห์เพื่อสร้างแผนผังภูมิทัศน์ทางวิชาการ (Academic Landscape)

\section{ปรัชญาการสืบค้นและคัดกรองวรรณกรรมวิชาการสากล}

การคัดเลือกวรรณกรรมที่มีคุณภาพ ต้องผ่านเกณฑ์การประเมิน 4 มิติ ได้แก่ ความสอดคล้องกับวัตถุประสงค์ (Relevance), ระเบียบวิธีวิจัยที่รัดกุม (Methodological Rigor), ความทันสมัยของผลการทดลอง (Recency) และความน่าเชื่อถือของแหล่งตีพิมพ์ (Journal Authority)

\begin{teachingplan}{การกำหนดคำค้นสำหรับการสืบค้นเชิงโครงสร้าง}
การสืบค้นผ่านฐานข้อมูล Scopus, Web of Science หรือ PubMed ควรกำหนดคำสำคัญแบบบูลีน (Boolean Search Operators) ร่วมกับการกำหนดกรอบปีตีพิมพ์ย้อนหลังไม่เกิน 5 ปี เพื่อให้ได้ข้อมูลที่สะท้อนแนวโน้มปัจจุบัน
\end{teachingplan}

\section{การสร้างเมทริกซ์เชิงธีม}

การจัดกลุ่มวรรณกรรมช่วยให้นักวิจัยมองเห็นรูปแบบแนวคิด (Thematic Patterns) ที่เกิดขึ้นซ้ำในชุมชนวิชาการ โดยสร้างตารางเปรียบเทียบดังนี้:

\begin{itemize}
    \item \textbf{แกนทฤษฎี (Theoretical Framework):} โมเดลทางฟิสิกส์ หรือสมการที่ใช้เป็นแกนหลัก
    \item \textbf{ตัวแปรที่ศึกษา (Variables):} ตัวแปรต้น ตัวแปรตาม และตัวแปรควบคุม
    \item \textbf{ระเบียบวิธีเก็บข้อมูล (Data Collection):} การทดลองจริง การใช้แบบจำลองคอมพิวเตอร์ หรือแบบวัดผลสัมฤทธิ์
    \item \textbf{ผลการค้นพบสำคัญ (Key Findings):} ค่าสถิติและข้อสรุปประจักษ์
\end{itemize}

\section{ระเบียบวิธีวิเคราะห์ช่องว่างงานวิจัย}

เมื่อจัดกลุ่มข้อมูลในเมทริกซ์แล้ว นักวิจัยจะสามารถจำแนกช่องว่างงานวิจัยได้ 3 ประเภทหลัก:

\begin{equation}
\text{Research Gap} = \text{Knowledge Domain} \setminus \bigcup_{i=1}^n \text{Reported Findings}_i
\end{equation}

ได้แก่ ช่องว่างด้านระเบียบวิธีวิจัย (Methodological Gap), ช่องว่างด้านกลุ่มตัวอย่างหรือบริบทประยุกต์ (Contextual Gap), และช่องว่างด้านผลลัพธ์ที่ขัดแย้งกัน (Contradictory Evidence Gap) ซึ่งเป็นฐานสำคัญในการตั้งสมมติฐานการวิจัยใหม่
